\documentclass{article}
\usepackage{graphicx}
\usepackage{url}

%% http://www.ctan.org/tex-archive/help/Catalogue/entries/igo.html
\usepackage{igo}
\gobansize{9}

\usepackage[empty,cm]{fullpage}
\begin{document}

%% The first page includes a 9x9 board.
\centerline{\Huge WPI Go Beginner's Night}
\centerline{\Huge 7pm Friday, February 27, 2009}
\centerline{\Huge @ Campus Center, 2nd floor}
\vfill

\begin{center}
\includegraphics{gb9.epsi}
\end{center}

\vfill

\newpage

\section*{Eye, eye, eye, eye, stayin' alive, stayin' alive...}

Go is often called the ``surrounding game'', and for good reason!  The
capturing rule applies when we surround stones.  But another thing we
can surround is empty space.  We may have noticed that groups that
surround a large amount of empty space are harder to capture.

In fact, if we surround in a particular way, our groups
become altogether immune.
\begin{center}
\shortstack{
\cleargoban
\black{a2,b2,b1,c2,d2,e2,e1}
\showgoban
}
\shortstack{
\cleargoban
\showgoban
}
\shortstack{
\cleargoban
\showgoban
}
%% diagrams here
\end{center}
The distinguishing feature about these invincible groups is that
they've surrounded two distinct empty areas.  Black can't play to
completely fill either area because suicidal moves are prohibited.  If
a group can form two eyes, even despite the opponent's best
resistance, then that group is \emph{alive}.


\section*{Puzzles}
Here are a few puzzles about the \emph{life and death} of groups.  Try to
solve them:
\begin{center}

%% put puzzles here on life and death

%% 

%% 
%%
%% 

\shortstack{
\showgoban\\
Black's move.\\Where can Black form eyes?
}
\hspace{.3in}
\shortstack{
\showgoban\\
Black's move.\\Where can Black form eyes?
}
\hspace{.3in}
\shortstack{
\showgoban\\
White's move.\\Can White interfere with Black's eyes?
}
\end{center}




\section*{Scoring}

Here is an example of a game at its end.
%% diagram here
The borders between Black and White are solid, and neither feels they
can make productive moves any more.  How do we score the game?
\begin{itemize}
\item Groups whose capture is inevitable are removed and given to the
  opponent.

\item Each player counts the whitespace sounded by their remaining
  groups.  This \emph{territory}, plus the number of captures, forms
  the player's score.
\end{itemize}




\section*{For more information...}
\begin{itemize}
\item \emph{WPI Go} (\url{http://go.hashcollision.org/}) is our homepage.

\item \emph{Sensei's Library} (\url{http://senseis.xmp.net/}), a wiki
  devoted to the game.

\item \emph{The Interactive Way to Go}
  (\url{http://playgo.to/interactive/}), which gives a more leisurely
  introduction.
\end{itemize}
\end{document}
